\begin{env}[1.2.4]
\label{0.1.2.4-fr}
El anillo de fracciones $S^{-1}A$ y el homomorfismo canónico $i_A^S$ son solución de un \emph{problema de aplicación universal}: todo homomorfismo $u$ de $A$ en un anillo $B$ tal que $u(S)$ está formado por elementos \emph{invertibles} de $B$ se factoriza de una única manera
\[
  u:A\xrightarrow[i_A^S]{}S^{-1}A\xrightarrow[u^*]{}B
\]

\oldpage[0\textsubscript{I}]{14}
donde $u^*$ es un homomorfismo de anillos.
Bajo las mismas hipótesis, sean $M$ un $A$-módulo, $N$ un $B$-módulo y $v:M\to N$ un homomorfismo de $A$-módulos (para la estructura de $B$-módulo en $N$ definida por $u:A\to B$); entonces $v$ se factoriza de una única manera
\[
  v:M\xrightarrow[i_M^S]{}S^{-1}M\xrightarrow[v^*]{}N
\]
donde $v^*$ es un homomorfismo de $S^{-1}A$-módulos (para la estructura de $S^{-1}A$-módulo en $N$ definida por $u^*$).
\end{env}

\begin{env}[1.2.5]
\label{0.1.2.5-fr}
Se define un isomorfismo canónico $S^{-1}A\otimes_A M\xrightarrow{\sim}S^{-1}M$ de $S^{-1}A$-módulos haciendo corresponder al elemento $(a/s)\otimes m$ el elemento $(am)/s$; el isomorfismo inverso envía $m/s$ a $(1/s)\otimes m$.
\end{env}

\begin{env}[1.2.6]
\label{0.1.2.6-fr}
Para todo ideal $\mathfrak{a}'$ de $S^{-1}A$, $\mathfrak{a}=(i_A^S)^{-1}(\mathfrak{a}')$ es un ideal de $A$, y $\mathfrak{a}'$ es el ideal de $S^{-1}A$ generado por $i_A^S(\mathfrak{a})$, que se identifica con $S^{-1}\mathfrak{a}$ (1.3.2).
La aplicación $\mathfrak{p}'\to(i_A^S)^{-1}(\mathfrak{p}')$ es un isomorfismo, para la estructura de orden, del conjunto de los ideales \emph{primos} de $S^{-1}A$ sobre el conjunto de los ideales primos $\mathfrak{p}$ de $A$ tales que $\mathfrak{p}\cap S=\varnothing$.
Además, los anillos locales $A_{\mathfrak{p}}$ y $(S^{-1}A)_{S^{-1}\mathfrak{p}}$ son entonces canónicamente (1.5.1) isomorfos.
\end{env}

\begin{env}[1.2.7]
\label{0.1.2.7-fr}
Cuando $A$ es un anillo \emph{íntegro}, cuyo cuerpo de fracciones se denota por $K$, la aplicación canónica $i_A^S:A\to S^{-1}A$ es inyectiva para todo subconjunto multiplicativo $S$ que no contiene $0$, y $S^{-1}A$ se identifica entonces canónicamente con un subanillo de $K$ que contiene $A$.
En particular, para todo ideal primo $\mathfrak{p}$ de $A$, $A_{\mathfrak{p}}$ es un anillo local que contiene $A$, de ideal maximal $\mathfrak{p}A_{\mathfrak{p}}$, y se tiene $\mathfrak{p}A_{\mathfrak{p}}\cap A=\mathfrak{p}$.
\end{env}

\begin{env}[1.2.8]
\label{0.1.2.8-fr}
Si $A$ es un anillo \emph{reducido} (1.1.1), lo mismo ocurre con $S^{-1}A$: en efecto, si $(x/s)^n=0$ para $x\in A$, $s\in S$, esto significa que existe $s'\in S$ tal que $s'x^n=0$, de donde $(s'x)^n=0$, lo que, por hipótesis, implica $s'x=0$, y por tanto $x/s=0$.
\end{env}
